\documentclass[12pt,a4paper]{article}

%\usepackage[backend=biber, style=numeric-comp, sorting=none]{biblatex}
%\addbibresource{refs.bib}

%\usepackage{graphicx, csquotes, tikz, pgfplots, wrapfig}
%\pgfplotsset{compat = newest}
%\graphicspath{{./pictures/}}

\usepackage[dutch]{babel}
\usepackage{csquotes}
\usepackage{amsmath, amssymb, amsthm}

\theoremstyle{definition}
\newtheorem*{notation*}{Notatie}

\usepackage{listings}

\lstdefinestyle{mystyle}{
    numbers=left,                    
    numbersep=5pt,                  
}
\lstset{style=mystyle}

\numberwithin{equation}{section}

\usepackage[hidelinks]{hyperref}

\binoppenalty=100000
\relpenalty=100000

\allowdisplaybreaks
\setlength\parindent{0pt}

\title{Verslag Taak 1}
\author{Robin Elskens en Lina Martens}
\date{} %TODO: Submission date

\begin{document}

\maketitle

In tegenstelling tot wat ik in de pastebin schreef is het best nuttig om na inductiehypothese, $a_i$ als de leidende coeëfficient van $P_i$ te zien om later wat schrijfwerk te besparen.
\subsubsection*{Graad van $P_n$}
We tonen de volgende formule aan met inductie: 

\begin{equation*}
\text{Graad$(P_n)=$}
\begin{cases}
n & \text{als $n$ even is}\\
n-1 & \text{als $n$ oneven is}
\end{cases}
%UNUSED: \label{eq:graad_P}. Remove * from equation when needed.
\end{equation*}
Zij $a_n$ de coëfficient van de term $r^g$ in $P_n$, waarbij $g=\text{Graad}(P_n)$. Als deze term niet bestaat zeggen we dat $a_n = 0$.
Merk op dat voor zowel $P_0$ als $P_1$ de graad 0 is, en dat dit in overeenstemming is met de formule.
Zij $k\in\mathbb{N}$, met $k\geq2$ en veronderstel dat de formule geldt voor elke $P_i$ met $i<k$. 
We merken op dat voor alle $i<k$ dan geldt dat $a_i$ de leidende coëfficient van $P_i$ is.
We hebben nu twee gevallen.\\
%
\underline{\textbf{1. $k$ is even}}\\
Uit de inductiehypothese volgt dat $\text{Graad}(P_{k-1}) = k-2$ en $\text{Graad}(P_{k-2}) = k-2$.
Volgens de recursieformule \eqref{eq:P_recursie} zal de graad van $P_k(r)$ dus bepaald worden door $-4r^2P_{k-2}(r)$. Er volgt dat $\text{Graad}(P_k) = k-2+2 = k$ en dat $a_n=-4a_{n-2}.$
\\
\underline{\textbf{2. $k$ is oneven}}\\
Uit de inductiehypothese volgt dat $\text{Graad}(P_{k-1}) = k-1$ en $\text{Graad}(P_{k-2}) = k-3$.
Volgens de recursieformule \eqref{eq:P_recursie} zal de graad van $P_k(r)$ dus bepaald worden door $(4k-2)P_{k-1}(r)-4r^2P_{k-2}(r)$.
Hieruit volgt dat de graad hoogstens $k-1$ zal zijn. Er geldt nu dat $a_{k} = (4k-2)a_{k-1} -4a_{k-2}$, waarbij we niet uitsluiten dat $a_k$ nul kan zijn.


\vspace{5mm}
Merk op dat we de formule dus bewezen hebben als we kunnen aantonen dat de rij $\left\{a_n\right\}_{n\in\mathbb{N}}$ nooit nul is. 
Inderdaad, uit puntje 1 en 2 in de inductie hierboven volgt dat de rij voldoet aan
$$a_n=
\begin{cases}
  -4 a_{n-2} & \text{als $n$ even is}\\
  \left(4n-2\right)a_{n-1} - 4a_{n-2} & \text{als $n$ oneven is}
\end{cases}$$
wanneer $n\geq2$.
Omdat het basisgeval $a_0 = 2$ niet nul is, volgt hieruit meteen dat geen enkel element met even index nul zal zijn.
Omdat het andere basisgeval $a_1 = 4$ is, is aan alle voorwaarden voor propositie \hyperref[thm:prop1]{1} voldaan.
Voor elke oneven $i$ geldt dus dat $$a_i = i(i+1)a_{i-1} \neq 0$$

%TODO: Analoog aan P_n
\subsubsection*{Graad van $Q_n$}
We tonen analoog als hierboven aan dat de volgende formule aan met inductie: 

\begin{equation*}
\text{Graad$(Q_n)=$}
\begin{cases}
  0 & \text{als $n=0$}\\
n-1 & \text{als $n$ even is}\\
n & \text{als $n$ oneven is}
\end{cases}
%UNUSED: \label{eq:graad_Q}. Remove * from equation when needed.
\end{equation*}
Merk op dat voor $Q_0(r) = 0$, $Q_1(r) = 4r$ en  $Q_2(r)=24r$ de formule geldt.
Zij $k\in\mathbb{N}$ en veronderstel dat dat de formule geldt voor elke $P_i$ met $i<k$. We hebben nu twee gevallen.\\
%
\underline{\textbf{1. $k$ is oneven}}\\
Uit de inductiehypothese volgt dat $\text{Graad}(Q_{k-1}) = k-2$ en $\text{Graad}(Q_{k-2}) = k-2$.
Volgens de recursieformule \eqref{eq:Q_recursie} zal de graad van $Q_k(r)$ dus bepaald worden door $-4r^2Q_{k-2}(r)$. Er volgt dat $\text{Graad}(Q_k) = k-2+2 = k.$
\\
\underline{\textbf{2. $k$ is even}}\\
Uit de inductiehypothese volgt dat $\text{Graad}(Q_{k-1}) = k-1$ en $\text{Graad}(Q_{k-2}) = k-3$.
Volgens de recursieformule \eqref{eq:Q_recursie} zal de graad van $Q_k(r)$ dus bepaald worden door $(4k-2)Q_{k-1}(r)-4r^2Q_{k-2}(r)$. We merken op dat de leidende coëfficienten van $a_i$ van $Q_i$ met $i<k$  niet nul zijn en analoog aan het bewijs voor de graad van $P_n$ voldoen aan de recursieformule uit propositie \hyperref[thm:prop2]{2}.
Sinds de leidende coëfficienten van $Q_0$, $Q_1$ en $Q_2$ resp. 0, 4 en 24 zijn, is aan  de voorwaarden van propositie \hyperref[thm:prop2]{2} voldaan, dus is $a_k = (4k-2)a_{k-1} - 4a_{k-2} = k(k+1)a_{k-1} \neq 0$.
Hieruit volgt dat $\text{Graad}(Q_k) = k-1$.

\end{document}
